\documentclass[12pt,a4paper]{article}
\usepackage[utf8]{inputenc}
\usepackage[T5]{fontenc}
\usepackage[vietnamese]{babel}
\usepackage{amsmath, amssymb}
\usepackage{array}
\usepackage{enumitem}
\usepackage{setspace}
\setstretch{1.18}
\usepackage{graphicx}
\usepackage{ifthen}
\usepackage{tikz}
\usetikzlibrary{arrows.meta, positioning, patterns}
\IfFileExists{pgf-pie.sty}{\usepackage{pgf-pie}}{}
\usepackage{pgfplots}
\pgfplotsset{compat=1.18}
\usepackage{chemfig}
\usepackage{longtable,booktabs}
\usepackage{adjustbox}
\providecommand{\errorpart}[2]{$\underset{\textbf{#1}}{\underline{\text{#2}}}$}
\usepackage[top=0.8in, left=0.75in, right=0.75in, bottom=0.85in, headheight=38pt, headsep=18pt, footskip=1.2cm]{geometry}
\usepackage{fancyhdr}
\usepackage[colorlinks=true, linkcolor=blue, urlcolor=blue]{hyperref}

\setlength{\parindent}{0pt}
\setlength{\parskip}{5pt plus 2pt minus 1pt}

\setlist[enumerate]{itemsep=4pt, parsep=2pt, topsep=4pt, partopsep=2pt}
\setlist[itemize]{itemsep=4pt, parsep=2pt, topsep=4pt}

\makeatletter
\AtEndDocument{\immediate\write\@auxout{\string\newlabel{LecturingLastPage}{{}{\thepage}{}{}{}}}}
\AtEndDocument{\immediate\write\@auxout{\string\newlabel{LecturingPageCount}{{}{\two@digits{\value{page}}}{}{}{}}}}
\makeatother
\pagestyle{fancy}
\fancyhf{}
\cfoot{\small Trang \thepage/\pageref*{LecturingLastPage} - Mã đề thi 9113}
\renewcommand{\headrulewidth}{0pt}
\renewcommand{\footrulewidth}{0pt}

\begin{document}

\sloppy

\noindent
\begin{minipage}[t]{0.37\textwidth}
    \small
    \centering
    \IfFileExists{logo.png}{\includegraphics[height=0.9cm]{logo.png}\\[0.1cm]}{%
      \IfFileExists{logo.jpg}{\includegraphics[height=0.9cm]{logo.jpg}\\[0.1cm]}{}%
    }
    {\bfseries \mbox{}}\\[-0.2cm]
    \rule{3cm}{0.5pt}\\[0.1cm]
    {\bfseries ĐỀ LUYỆN TẬP}\\{\small (Đề thi có \pageref*{LecturingPageCount} trang)}
\end{minipage}%
\hfill
\begin{minipage}[t]{0.62\textwidth}
    \small
    \centering
    {\bfseries \sbox0{KỲ THI TỐT NGHIỆP TRUNG HỌC PHỔ THÔNG NĂM 2026}\ifdim\wd0>\linewidth KỲ THI TỐT NGHIỆP TRUNG HỌC PHỔ~THÔNG\\NĂM 2026\else\usebox0\fi}\\[0.15cm]
    {\bfseries Môn thi: VẬT LÍ}\\[0.15cm]
    {\itshape Thời gian làm bài: 50 phút, không kể thời gian phát đề}
\end{minipage}

\vspace{0.4cm}
\noindent
\begin{minipage}[t]{0.68\textwidth}
    Họ, tên thí sinh: \dotfill\\[0.15cm]
    Số báo danh: \dotfill
\end{minipage}%
\hfill
\fbox{\begin{minipage}{3.6cm}\centering\textbf{Mã đề: 9113}\end{minipage}}

\vspace{0.5cm}

% LECTURING-PART: P1
\textbf{Phần I: Trắc nghiệm nhiều lựa chọn} \textit{Thí sinh trả lời từ câu 1 đến câu 18. Mỗi câu hỏi thí sinh chỉ chọn một phương án.}

\noindent + Cho biết: $\pi = 3{,}14$; $T(\text{K}) = t(^{\circ}\text{C}) + 273$; $R = 8{,}31\ \text{J.mol}^{-1}.\text{K}^{-1}$; $N_A = 6{,}02 \cdot 10^{23}$ hạt/mol; $1\ \text{u} = 931{,}5\ \text{MeV}/c^{2}$; $\mu_0 = 4\pi \cdot 10^{-7}\ \text{T.m/A}$.\\
\noindent + Không làm tròn kết quả các phép tính trung gian.\\

\begin{minipage}{\textwidth}
\noindent \textbf{Câu 1:} Một lượng khí lí tưởng xác định được nung nóng đẳng tích từ $300\ \text{K}$ đến $450\ \text{K}$. Biết áp suất ban đầu của khí là $2{,}0\ \text{atm}$. Áp suất của khí ở trạng thái sau là\\
\noindent\makebox[0.48\textwidth][l]{A. $1{,}3\ \text{atm}$.}\makebox[0.48\textwidth][l]{B. $2{,}5\ \text{atm}$.}\\
\noindent\makebox[0.48\textwidth][l]{C. $3{,}0\ \text{atm}$.}\makebox[0.48\textwidth][l]{D. $4{,}5\ \text{atm}$.}
\end{minipage}


\begin{minipage}{\textwidth}
\noindent \textbf{Câu 2:} Một khối khí nhận nhiệt lượng $120\ \text{J}$ từ môi trường và đồng thời thực hiện công $50\ \text{J}$ lên môi trường. Độ biến thiên nội năng của khối khí là\\
\noindent\makebox[0.48\textwidth][l]{A. $+70\ \text{J}$.}\makebox[0.48\textwidth][l]{B. $-70\ \text{J}$.}\\
\noindent\makebox[0.48\textwidth][l]{C. $-170\ \text{J}$.}\makebox[0.48\textwidth][l]{D. $+170\ \text{J}$.}
\end{minipage}


\begin{minipage}{\textwidth}
\noindent \textbf{Câu 3:} Truyền nhiệt lượng $84\ \text{kJ}$ cho $2{,}0\ \text{kg}$ nước, biết nhiệt dung riêng của nước là $4200\ \text{J/(kg.K)}$ và không có hao phí nhiệt ra môi trường. Nhiệt độ của nước tăng thêm\\
\noindent A. $10\ ^{\circ}\text{C}$. \hfill B. $20\ ^{\circ}\text{C}$. \hfill C. $5\ ^{\circ}\text{C}$. \hfill D. $40\ ^{\circ}\text{C}$.
\end{minipage}


\begin{minipage}{\textwidth}
\noindent \textbf{Câu 4:} Khi nhiệt độ tuyệt đối của một khối khí lí tưởng tăng lên gấp đôi thì động năng tịnh tiến trung bình của các phân tử khí\\
\noindent\makebox[0.48\textwidth][l]{A. tăng gấp $\sqrt{2}$ lần.}\makebox[0.48\textwidth][l]{B. không đổi.}\\
\noindent\makebox[0.48\textwidth][l]{C. tăng gấp $2$ lần.}\makebox[0.48\textwidth][l]{D. tăng gấp $4$ lần.}
\end{minipage}


\begin{minipage}{\textwidth}
\noindent \textbf{Câu 5:} Hình bên là đồ thị biểu diễn sự phụ thuộc của thể tích $V$ vào nhiệt độ tuyệt đối $T$ của một lượng khí lí tưởng xác định trong một quá trình biến đổi trạng thái. Đồ thị là đoạn thẳng có đường kéo dài đi qua gốc tọa độ. Khi nhiệt độ của khí là $450\ \text{K}$ thì thể tích của khí là\\
\begin{center}
\begin{adjustbox}{max width=0.95\linewidth, max totalheight=0.38\textheight}
\begin{tikzpicture}[>=stealth, thick, every node/.style={font=\small}]
  % trục và đồ thị
  \draw[->] (0,0) -- (6.5,0);
  \draw[->] (0,0) -- (0,5.6);
  \draw[thick] (0,0) -- (3.6,3);
  \draw[dashed] (3.6,0) -- (3.6,3);
  \draw[dashed] (0,3) -- (3.6,3);
  \fill (3.6,3) circle (1.5pt);
  \fill (0,0) circle (1.2pt);
  % nhãn
    % captions last so nothing paints over them
    % captions last so nothing paints over them
    \node[below left] at (0,0) {$O$};
\node[right] at (6.5,0) {$T\ (\text{K})$};
\node[above] at (0,5.6) {$V\ (\text{lít})$};
\node[below] at (3.6,-0.05) {$300$};
\node[left] at (-0.05,3) {$6$};
\node[above left, fill=white, inner sep=2pt] at (3.6,3) {$M$};
\end{tikzpicture}
\end{adjustbox}
\end{center}
\noindent\makebox[0.24\textwidth][l]{A. $4$ lít.}\makebox[0.24\textwidth][l]{B. $7{,}5$ lít.}\makebox[0.24\textwidth][l]{C. $9$ lít.}\makebox[0.24\textwidth][l]{D. $12$ lít.}
\end{minipage}


\begin{minipage}{\textwidth}
\noindent \textbf{Câu 6:} Một đoạn dây dẫn thẳng dài $0{,}5\ \text{m}$ có dòng điện cường độ $4\ \text{A}$ chạy qua, đặt trong từ trường đều có cảm ứng từ $0{,}2\ \text{T}$ sao cho dây hợp với các đường sức từ một góc $30^{\circ}$. Độ lớn lực từ tác dụng lên đoạn dây là\\
\noindent\makebox[0.48\textwidth][l]{A. $0{,}1\ \text{N}$.}\makebox[0.48\textwidth][l]{B. $0{,}2\ \text{N}$.}\\
\noindent\makebox[0.48\textwidth][l]{C. $0{,}35\ \text{N}$.}\makebox[0.48\textwidth][l]{D. $0{,}4\ \text{N}$.}
\end{minipage}


\begin{minipage}{\textwidth}
\noindent \textbf{Câu 7:} Một khung dây phẳng có diện tích $100\ \text{cm}^2$ đặt trong từ trường đều có cảm ứng từ $0{,}5\ \text{T}$, mặt phẳng khung dây vuông góc với các đường sức từ. Từ thông qua khung dây là\\
\noindent\makebox[0.48\textwidth][l]{A. $50\ \text{Wb}$.}\makebox[0.48\textwidth][l]{B. $5 \cdot 10^{-3}\ \text{Wb}$.}\\
\noindent\makebox[0.48\textwidth][l]{C. $5 \cdot 10^{-2}\ \text{Wb}$.}\makebox[0.48\textwidth][l]{D. $0{,}5\ \text{Wb}$.}
\end{minipage}


\begin{minipage}{\textwidth}
\noindent \textbf{Câu 8:} Từ thông qua một cuộn dây gồm $200$ vòng giảm đều từ $0{,}06\ \text{Wb}$ xuống $0{,}02\ \text{Wb}$ trong khoảng thời gian $0{,}1\ \text{s}$. Độ lớn suất điện động cảm ứng xuất hiện trong cuộn dây là\\
\noindent\makebox[0.48\textwidth][l]{A. $0{,}4\ \text{V}$.}\makebox[0.48\textwidth][l]{B. $8\ \text{V}$.}\\
\noindent\makebox[0.48\textwidth][l]{C. $40\ \text{V}$.}\makebox[0.48\textwidth][l]{D. $80\ \text{V}$.}
\end{minipage}


\begin{minipage}{\textwidth}
\noindent \textbf{Câu 9:} Cho mạch điện như hình bên. Nguồn điện có suất điện động $12\ \text{V}$, điện trở trong $2\ \Omega$; các điện trở $R_1 = 4\ \Omega$, $R_2 = 6\ \Omega$, $R_3 = 3\ \Omega$; ampe kế $A$ lí tưởng mắc nối tiếp với $R_2$. Số chỉ của ampe kế là\\
\begin{center}
\begin{adjustbox}{max width=0.95\linewidth, max totalheight=0.38\textheight}
\begin{tikzpicture}[thick, every node/.style={font=\small}]
  % dây dẫn
  \draw (0,1.7) -- (0,3) -- (1.1,3);
  \draw (2.3,3) -- (3.2,3);
  \draw (3.2,3) -- (3.9,3);
  \draw (4.5,3) -- (5.2,3);
  \draw (6.4,3) -- (7,3) -- (7,0);
  \draw (3.2,3) -- (3.2,2.2);
  \draw (3.2,1.2) -- (3.2,0);
  \draw (7,0) -- (0,0) -- (0,1.3);
  % nguồn: bản dài là cực dương (phía trên), bản ngắn là cực âm
  \draw (-0.4,1.7) -- (0.4,1.7);
  \draw[very thick] (-0.2,1.3) -- (0.2,1.3);
  % điện trở R1, R2, R3
  \draw[fill=white] (1.1,2.8) rectangle (2.3,3.2);
  \draw[fill=white] (5.2,2.8) rectangle (6.4,3.2);
  \draw[fill=white] (3.0,1.2) rectangle (3.4,2.2);
  % ampe kế
  \draw[fill=white] (4.2,3) circle (0.3);
  % nút nối
  \fill (3.2,3) circle (2pt);
  \fill (3.2,0) circle (2pt);
  % nhãn đặt sau cùng
    % captions last so nothing paints over them
    % captions last so nothing paints over them
    \node[fill=white, inner sep=1.5pt] at (4.2,3) {$A$};
\node[above] at (1.7,3.2) {$R_1 = 4\ \Omega$};
\node[above] at (5.8,3.2) {$R_2 = 6\ \Omega$};
\node[right=3pt] at (3.4,1.7) {$R_3 = 3\ \Omega$};
\node[left=6pt] at (0,1.5) {$E,\ r$};
\node[above right] at (0.1,1.7) {$+$};
\node[below right] at (0.1,1.3) {$-$};
\end{tikzpicture}
\end{adjustbox}
\end{center}
\noindent\makebox[0.24\textwidth][l]{A. $1{,}5\ \text{A}$.}\makebox[0.24\textwidth][l]{B. $1{,}0\ \text{A}$.}\makebox[0.24\textwidth][l]{C. $0{,}25\ \text{A}$.}\makebox[0.24\textwidth][l]{D. $0{,}5\ \text{A}$.}
\end{minipage}


\begin{minipage}{\textwidth}
\noindent \textbf{Câu 10:} Độ hụt khối của hạt nhân $^{4}_{2}\text{He}$ là $0{,}0304\ \text{u}$. Năng lượng liên kết riêng của hạt nhân này xấp xỉ bằng\\
\noindent\makebox[0.48\textwidth][l]{A. $3{,}54\ \text{MeV/nuclôn}$.}\makebox[0.48\textwidth][l]{B. $28{,}3\ \text{MeV/nuclôn}$.}\\
\noindent\makebox[0.48\textwidth][l]{C. $14{,}2\ \text{MeV/nuclôn}$.}\makebox[0.48\textwidth][l]{D. $7{,}08\ \text{MeV/nuclôn}$.}
\end{minipage}


\begin{minipage}{\textwidth}
\noindent \textbf{Câu 11:} Một chất phóng xạ có chu kì bán rã $8$ ngày. Sau $24$ ngày, tỉ lệ phần trăm số hạt nhân phóng xạ còn lại so với ban đầu là\\
\noindent\makebox[0.24\textwidth][l]{A. $12{,}5\%$.}\makebox[0.24\textwidth][l]{B. $6{,}25\%$.}\makebox[0.24\textwidth][l]{C. $25\%$.}\makebox[0.24\textwidth][l]{D. $37{,}5\%$.}
\end{minipage}


\begin{minipage}{\textwidth}
\noindent \textbf{Câu 12:} Cho phản ứng hạt nhân $^{14}_{7}\text{N} + {}^{4}_{2}\text{He} \to {}^{17}_{8}\text{O} + X$. Hạt $X$ là\\
\noindent\makebox[0.48\textwidth][l]{A. prôtôn $^{1}_{1}\text{H}$.}\makebox[0.48\textwidth][l]{B. êlectron $^{0}_{-1}\text{e}$.}\\
\noindent\makebox[0.48\textwidth][l]{C. đơteri $^{2}_{1}\text{H}$.}\makebox[0.48\textwidth][l]{D. nơtron $^{1}_{0}\text{n}$.}
\end{minipage}


\begin{minipage}{\textwidth}
\noindent \textbf{Câu 13:} Một con lắc lò xo gồm vật nhỏ khối lượng $0{,}25\ \text{kg}$ gắn với lò xo có độ cứng $100\ \text{N/m}$, dao động điều hòa. Chu kì dao động của con lắc xấp xỉ bằng\\
\noindent\makebox[0.48\textwidth][l]{A. $0{,}16\ \text{s}$.}\makebox[0.48\textwidth][l]{B. $0{,}31\ \text{s}$.}\\
\noindent\makebox[0.48\textwidth][l]{C. $0{,}63\ \text{s}$.}\makebox[0.48\textwidth][l]{D. $3{,}14\ \text{s}$.}
\end{minipage}


\begin{minipage}{\textwidth}
\noindent \textbf{Câu 14:} Một lượng khí lí tưởng ở $27\ ^{\circ}\text{C}$ có áp suất $2{,}0\ \text{atm}$ và thể tích $4{,}0$ lít. Nung nóng khí đến $127\ ^{\circ}\text{C}$ thì thể tích của khí là $5{,}0$ lít. Áp suất của khí lúc này xấp xỉ bằng\\
\noindent\makebox[0.48\textwidth][l]{A. $1{,}5\ \text{atm}$.}\makebox[0.48\textwidth][l]{B. $2{,}13\ \text{atm}$.}\\
\noindent\makebox[0.48\textwidth][l]{C. $2{,}67\ \text{atm}$.}\makebox[0.48\textwidth][l]{D. $3{,}0\ \text{atm}$.}
\end{minipage}


\begin{minipage}{\textwidth}
\noindent \textbf{Câu 15:} Một dây dẫn thẳng rất dài có dòng điện cường độ $5\ \text{A}$ chạy qua, đặt trong không khí. Độ lớn cảm ứng từ do dòng điện gây ra tại điểm cách dây $10\ \text{cm}$ là\\
\noindent\makebox[0.48\textwidth][l]{A. $5 \cdot 10^{-6}\ \text{T}$.}\makebox[0.48\textwidth][l]{B. $10^{-5}\ \text{T}$.}\\
\noindent\makebox[0.48\textwidth][l]{C. $2 \cdot 10^{-5}\ \text{T}$.}\makebox[0.48\textwidth][l]{D. $10^{-6}\ \text{T}$.}
\end{minipage}


\begin{minipage}{\textwidth}
\noindent \textbf{Câu 16:} Điện áp xoay chiều giữa hai đầu một đoạn mạch có biểu thức $u = 220\sqrt{2}\cos(100\pi t)\ \text{V}$. Điện áp hiệu dụng và tần số của dòng điện lần lượt là\\
\noindent A. $220\ \text{V}$; $50\ \text{Hz}$.\\
\noindent B. $220\ \text{V}$; $100\ \text{Hz}$.\\
\noindent C. $311\ \text{V}$; $100\ \text{Hz}$.\\
\noindent D. $220\sqrt{2}\ \text{V}$; $50\ \text{Hz}$.
\end{minipage}


\begin{minipage}{\textwidth}
\noindent \textbf{Câu 17:} Một máy biến áp lí tưởng có cuộn sơ cấp $1000$ vòng, cuộn thứ cấp $50$ vòng. Đặt vào hai đầu cuộn sơ cấp điện áp xoay chiều có giá trị hiệu dụng $220\ \text{V}$. Điện áp hiệu dụng ở hai đầu cuộn thứ cấp để hở là\\
\noindent\makebox[0.48\textwidth][l]{A. $11\ \text{V}$.}\makebox[0.48\textwidth][l]{B. $22\ \text{V}$.}\\
\noindent\makebox[0.48\textwidth][l]{C. $110\ \text{V}$.}\makebox[0.48\textwidth][l]{D. $4400\ \text{V}$.}
\end{minipage}


\begin{minipage}{\textwidth}
\noindent \textbf{Câu 18:} Nhiệt hóa hơi riêng của nước ở $100\ ^{\circ}\text{C}$ là $2{,}3 \cdot 10^{6}\ \text{J/kg}$. Nhiệt lượng cần cung cấp để $0{,}5\ \text{kg}$ nước ở $100\ ^{\circ}\text{C}$ hóa hơi hoàn toàn ở nhiệt độ đó là\\
\noindent\makebox[0.48\textwidth][l]{A. $4{,}6 \cdot 10^{6}\ \text{J}$.}\makebox[0.48\textwidth][l]{B. $1{,}15 \cdot 10^{6}\ \text{J}$.}\\
\noindent\makebox[0.48\textwidth][l]{C. $1{,}15 \cdot 10^{5}\ \text{J}$.}\makebox[0.48\textwidth][l]{D. $2{,}3 \cdot 10^{6}\ \text{J}$.}
\end{minipage}

\vspace{0.5cm}
% LECTURING-PART: P2
\textbf{Phần II: Trắc nghiệm đúng/sai} \textit{Thí sinh trả lời từ câu 1 đến câu 4. Trong mỗi ý a), b), c), d) ở mỗi câu, thí sinh chọn đúng hoặc sai.}

\begin{minipage}{\textwidth}
\noindent \textbf{Câu 1:} Một nhóm học sinh xác định hiệu suất của một dây đun điện nhúng trong nhiệt lượng kế như hình bên. Nhiệt lượng kế chứa $0{,}5\ \text{kg}$ nước, dây đun có công suất không đổi $100\ \text{W}$. Nhiệt dung riêng của nước là $4200\ \text{J/(kg.K)}$. Sau $210\ \text{s}$ đun, nhiệt độ của nước tăng từ $20\ ^{\circ}\text{C}$ lên $28\ ^{\circ}\text{C}$.\\
\begin{center}
\begin{adjustbox}{max width=0.95\linewidth, max totalheight=0.38\textheight}
\begin{tikzpicture}[thick, every node/.style={font=\small}]
  % nước
  \fill[blue!10] (0.05,0.05) rectangle (4.35,2.5);
  % cốc nhiệt lượng kế
  \draw (0,3.3) -- (0,0) -- (4.4,0) -- (4.4,3.3);
  % nhiệt kế
  \fill[white] (0.5,1.0) rectangle (0.7,3.9);
  \draw (0.5,1.0) rectangle (0.7,3.9);
  \fill[gray!60] (0.6,0.85) circle (0.17);
  \draw (0.6,0.85) circle (0.17);
  % que khuấy
  \draw[very thick] (1.4,0.4) -- (1.4,3.9);
  \draw (1.15,0.4) -- (1.65,0.4);
  % dây đun hình lò xo
  \draw (2.0,0.7) -- (2.1,0.95) -- (2.3,0.45) -- (2.5,0.95) -- (2.7,0.45) -- (2.9,0.95) -- (3.1,0.45) -- (3.3,0.95) -- (3.4,0.7);
  \draw (2.0,0.7) -- (2.0,4.2) -- (6,4.2) -- (6,4.05);
  \draw (3.4,0.7) -- (3.4,3.7) -- (6,3.7) -- (6,3.85);
  % nguồn điện
  \draw (5.7,4.05) -- (6.3,4.05);
  \draw[very thick] (5.85,3.85) -- (6.15,3.85);
  % nhãn đặt sau cùng

  \draw[thin] (-0.45,3.0) -- (0.5,3.0);

  \draw[thin] (0.9,4.4) -- (1.4,3.9);

  \draw[thin] (4.95,1.4) -- (3.1,0.95);

  \draw[thin] (4.95,0.5) -- (4.2,0.5);
    % captions last so nothing paints over them
    % captions last so nothing paints over them
    \node[anchor=east] at (-0.5,3.0) {nhiệt kế};
\node[anchor=south] at (0.5,4.45) {que khuấy};
\node[anchor=west] at (5.0,1.4) {dây đun};
\node[anchor=west] at (5.0,0.5) {nước};
\node[right=4pt] at (6.3,3.95) {nguồn};
\end{tikzpicture}
\end{adjustbox}
\end{center}
\noindent a) Nhiệt lượng do dây đun tỏa ra trong $210\ \text{s}$ là $21\ \text{kJ}$.\\

\noindent b) Nhiệt kế đặt sát dây đun cho số chỉ phản ánh chính xác nhất nhiệt độ trung bình của nước trong cốc.\\

\noindent c) Nhiệt lượng mà nước thu vào trong $210\ \text{s}$ là $16{,}8\ \text{kJ}$.\\

\noindent d) Hiệu suất của quá trình đun nước trong thí nghiệm này bằng $90\%$.
\end{minipage}


\begin{minipage}{\textwidth}
\noindent \textbf{Câu 2:} Một xilanh đặt thẳng đứng, miệng hở hướng lên trên, bên trong chứa một lượng khí lí tưởng được ngăn với bên ngoài bằng pít-tông nhẹ, diện tích $20\ \text{cm}^2$, có thể trượt không ma sát. Áp suất khí quyển là $1{,}0 \cdot 10^{5}\ \text{Pa}$. Ban đầu khí ở nhiệt độ $27\ ^{\circ}\text{C}$ và cột khí cao $30\ \text{cm}$. Khí được nung nóng chậm đến $87\ ^{\circ}\text{C}$, pít-tông dịch chuyển lên trên.\\
\noindent a) Trong quá trình nung, áp suất của khí trong xilanh luôn bằng áp suất khí quyển.\\

\noindent b) Trong quá trình này, thể tích của khí tỉ lệ thuận với nhiệt độ của khí tính theo thang Celsius.\\

\noindent c) Khi nhiệt độ của khí là $87\ ^{\circ}\text{C}$, cột khí cao $36\ \text{cm}$.\\

\noindent d) Công mà khí thực hiện lên môi trường ngoài trong quá trình nung là $12\ \text{J}$.
\end{minipage}


\begin{minipage}{\textwidth}
\noindent \textbf{Câu 3:} Một thanh dẫn $MN$ dài $0{,}4\ \text{m}$ trượt đều với tốc độ $5\ \text{m/s}$ trên hai thanh ray song song nằm ngang, trong một từ trường đều có cảm ứng từ $0{,}5\ \text{T}$ hướng thẳng đứng từ dưới lên như hình bên. Hai đầu ray nối với điện trở $R = 2\ \Omega$; điện trở của thanh và của ray không đáng kể.\\
\begin{center}
\begin{adjustbox}{max width=0.95\linewidth, max totalheight=0.38\textheight}
\begin{tikzpicture}[>=stealth, thick, every node/.style={font=\small}]
  % ray và điện trở
  \draw (0.8,0) -- (8,0);
  \draw (0.8,3) -- (8,3);
  \draw (0.8,0) -- (0.8,1.1);
  \draw (0.8,1.9) -- (0.8,3);
  \draw[fill=white] (0.6,1.1) rectangle (1.0,1.9);
  % từ trường hướng ra ngoài mặt giấy
  \foreach \x in {2.3,3.4,6.9,7.6} {
    \foreach \y in {0.7,2.3} {
      \draw[thin] (\x,\y) circle (0.1);
      \fill (\x,\y) circle (0.03);
    }
  }
  % thanh MN
  \draw[line width=2pt] (4.5,-0.15) -- (4.5,3.15);
  % vận tốc
  \draw[->] (4.8,1.5) -- (6.0,1.5);
  % nhãn đặt sau cùng
    % captions last so nothing paints over them
    % captions last so nothing paints over them
    \node[right=4pt] at (1.0,1.5) {$R$};
\node[above] at (4.5,3.15) {$M$};
\node[below] at (4.5,-0.15) {$N$};
\node[above] at (5.4,1.5) {$\overrightarrow{v}$};
\node[fill=white, inner sep=1.5pt] at (2.85,1.5) {$\overrightarrow{B}$};
\end{tikzpicture}
\end{adjustbox}
\end{center}
\noindent a) Suất điện động cảm ứng xuất hiện trong thanh $MN$ có độ lớn $1\ \text{V}$.\\

\noindent b) Cường độ dòng điện chạy trong mạch là $0{,}5\ \text{A}$.\\

\noindent c) Lực từ tác dụng lên thanh $MN$ cùng chiều với chiều chuyển động của thanh.\\

\noindent d) Công suất cơ học cần cung cấp để thanh chuyển động đều bằng $0{,}5\ \text{W}$.
\end{minipage}


\begin{minipage}{\textwidth}
\noindent \textbf{Câu 4:} Một nhóm học sinh đo chu kì bán rã của một mẫu phóng xạ bằng ống đếm Geiger. Trước tiên, nhóm đo khi không có mẫu và ghi được số xung nền ổn định là $20$ xung/phút. Sau đó, nhóm đặt mẫu trước ống đếm và đếm số xung trong $1$ phút tại các thời điểm khác nhau, kết quả như sau:\\
\begin{center}
\begin{adjustbox}{max width=\linewidth}
\begin{tabular}{|c|c|c|c|}
\hline
Thời điểm $t$ (phút) & $0$ & $10$ & $20$\\
\hline
Số xung đếm được trong $1$ phút & $340$ & $180$ & $100$\\
\hline
\end{tabular}
\end{adjustbox}
\end{center}
\noindent a) Tại thời điểm $t = 0$, số xung/phút do mẫu phóng xạ gây ra là $320$.\\

\noindent b) Chu kì bán rã của mẫu phóng xạ này là $20$ phút.\\

\noindent c) Nếu nhóm không trừ số xung nền mà dùng trực tiếp số liệu thô tại $t = 0$ và $t = 10$ phút để tính chu kì bán rã thì kết quả thu được nhỏ hơn giá trị thực.\\

\noindent d) Tại thời điểm $t = 30$ phút, số xung đếm được trong $1$ phút (đã gồm cả xung nền) dự kiến là $60$.
\end{minipage}

\vspace{0.5cm}
% LECTURING-PART: P3
\textbf{Phần III: Trắc nghiệm trả lời ngắn} \textit{Thí sinh trả lời từ câu 1 đến câu 6. Thí sinh chỉ điền một con số (có thể có dấu âm, dấu phẩy thập phân) vào ô trống.}

\noindent \textbf{Dùng thông tin sau cho Câu 1 và Câu 2:} Một lốp xe ô tô có thể tích không đổi $30$ lít, chứa không khí coi là khí lí tưởng. Khi xe đỗ trong bóng râm, nhiệt độ không khí trong lốp là $27\ ^{\circ}\text{C}$ và áp suất là $2{,}2\ \text{atm}$. Sau khi xe chạy một quãng đường dài, nhiệt độ không khí trong lốp tăng lên $57\ ^{\circ}\text{C}$.\\\par\nopagebreak[4]

\begin{minipage}{\textwidth}
\noindent \textbf{Câu 1:} Áp suất của không khí trong lốp sau khi xe chạy là bao nhiêu atm? \textit{(Làm tròn đến chữ số hàng phần mười)}
\end{minipage}


\begin{minipage}{\textwidth}
\noindent \textbf{Câu 2:} Người lái xe mở van xả bớt một phần không khí để áp suất trong lốp trở về $2{,}2\ \text{atm}$ trong khi nhiệt độ vẫn là $57\ ^{\circ}\text{C}$. Tìm $x$ biết khối lượng không khí đã thoát ra bằng $x\%$ khối lượng không khí trong lốp ngay trước khi xả. \textit{(Làm tròn đến chữ số hàng phần mười)}\\

\noindent \textbf{Dùng thông tin sau cho Câu 3 và Câu 4:} Một mẫu chất phóng xạ có chu kì bán rã $T = 12$ giờ. Tại thời điểm ban đầu, độ phóng xạ của mẫu là $H_0 = 4{,}0 \cdot 10^{6}\ \text{Bq}$. Lấy $\ln 2 = 0{,}693$ và $1$ giờ $= 3600\ \text{s}$.\\\par\nopagebreak[4]
\end{minipage}


\begin{minipage}{\textwidth}
\noindent \textbf{Câu 3:} Sau $30$ giờ kể từ thời điểm ban đầu, độ phóng xạ của mẫu là $x \cdot 10^{5}\ \text{Bq}$. Tìm $x$. \textit{(Làm tròn đến chữ số hàng phần trăm)}
\end{minipage}


\begin{minipage}{\textwidth}
\noindent \textbf{Câu 4:} Số hạt nhân phóng xạ có trong mẫu tại thời điểm ban đầu là $x \cdot 10^{11}$ hạt. Tìm $x$. \textit{(Làm tròn đến chữ số hàng phần trăm)}\\

\noindent \textbf{Dùng thông tin sau cho Câu 5 và Câu 6:} Một máy phát điện xoay chiều một pha đơn giản gồm khung dây dẫn có $100$ vòng, diện tích mỗi vòng $200\ \text{cm}^2$, quay đều quanh một trục nằm trong mặt phẳng khung và vuông góc với các đường sức của từ trường đều có cảm ứng từ $0{,}06\ \text{T}$, với tốc độ góc $100\ \text{rad/s}$. Khung dây có điện trở trong $2\ \Omega$.\\\par\nopagebreak[4]
\end{minipage}


\begin{minipage}{\textwidth}
\noindent \textbf{Câu 5:} Giá trị hiệu dụng của suất điện động xuất hiện trong khung dây là bao nhiêu vôn? \textit{(Làm tròn đến chữ số hàng phần mười)}
\end{minipage}


\begin{minipage}{\textwidth}
\noindent \textbf{Câu 6:} Nối hai đầu khung dây với một điện trở thuần $R = 3\ \Omega$ thành mạch kín. Công suất tỏa nhiệt trung bình trên điện trở $R$ là bao nhiêu oát? \textit{(Làm tròn đến chữ số hàng phần mười)}\\\vspace{0.6cm}
\begin{center}
{\bfseries -{}-{}-{}-{}-{}-}\ \textbf{HẾT}\ {\bfseries -{}-{}-{}-{}-{}-}
\end{center}
\vspace{0.15cm}
{\small\itshape
Thí sinh không được sử dụng tài liệu. Cán bộ coi thi không giải thích gì thêm.}

\end{minipage}
\end{document}
